% Severity-Aware Conformal Risk Control (Section 4) -- ML4H/PMLR Overleaf project.
% Requires: amsmath, amssymb, booktabs (worked-example table).
% For the Proposition/proof, add to your preamble:
%     \usepackage{amsthm}
%     \newtheorem{proposition}{Proposition}
% Depends on notation and equations from Section 3 (Preliminaries): \eqref{eq:loss},
% \eqref{eq:emp-risk}, \eqref{eq:crc-threshold}, \eqref{eq:crc-guarantee}.
% Bib keys used: angelopoulos2024crc, vovk2012conditional.

\section{Severity-Aware Conformal Risk Control}
\label{sec:sac}

The marginal guarantee in \eqref{eq:crc-guarantee} averages over tiers, which is
precisely its weakness. Because \textsc{benign} claims are far more numerous than
\textsc{dangerous} ones, the mean loss can sit below $\alpha$ while the loss
restricted to \textsc{dangerous} claims runs well above it: the abundant, easily
filtered benign claims subsidize the dangerous ones. A single budget cannot see
this, because it only ever constrains the pooled rate.

Severity-aware conformal risk control (SAC) removes the pooling. We partition the
calibration set by tier, $\mathcal{C}_t = \{\, i : t(c_i) = t \,\}$ with
$n_t = |\mathcal{C}_t|$, assign each tier its own budget $\alpha_t$, and run the
conformal risk control procedure of
\eqref{eq:emp-risk}--\eqref{eq:crc-threshold} \emph{independently within each tier}:
\begin{align}
\widehat{R}_{n_t}^{\,t}(\lambda) &= \frac{1}{n_t}\sum_{i \in \mathcal{C}_t}
\ell(c_i; \lambda), \label{eq:sac-risk}\\
\widehat{\lambda}_t &= \inf\left\{ \lambda :\;
\frac{n_t}{n_t+1}\,\widehat{R}_{n_t}^{\,t}(\lambda) + \frac{B}{n_t+1}
\;\le\; \alpha_t \right\}. \label{eq:sac-threshold}
\end{align}
At deployment, a claim is kept when its score clears the threshold \emph{of its own
tier}: keep $c$ iff $s(c) \ge \widehat{\lambda}_{t(c)}$. In our experiments we set a
strict dangerous budget and a loose benign budget, $\alpha_{\textsc{dangerous}} =
0.05$ and $\alpha_{\textsc{benign}} = 0.15$, against a marginal baseline of
$\alpha = 0.10$. This is the Mondrian, or group-conditional, construction
\citep{vovk2012conditional} applied to a clinical severity taxonomy: conditioning
the guarantee on the tier rather than reweighting claims inside one budget.

The per-tier guarantee follows immediately, which is the point.

\begin{proposition}[Per-tier risk control]
\label{prop:sac}
Fix a tier $t \in \mathcal{T}$ and suppose the calibration claims
$\{c_i : i \in \mathcal{C}_t\}$ and a fresh test claim of tier $t$ are
exchangeable. Then the tier-$t$ threshold $\widehat{\lambda}_t$ from
\eqref{eq:sac-threshold} satisfies
\begin{equation}
\mathbb{E}\!\left[\ell\big(c; \widehat{\lambda}_t\big) \,\middle|\, t(c) = t\right]
\;\le\; \alpha_t .
\end{equation}
Consequently the bounds hold simultaneously across tiers: the expected rate at
which a \emph{dangerous} hallucination is kept is at most
$\alpha_{\textsc{dangerous}}$, regardless of the tiers' base rates or their
relative sizes.
\end{proposition}

\begin{proof}
Restrict attention to tier $t$. On the subsequence $\mathcal{C}_t$ together with a
test claim of the same tier, the loss $\ell(\cdot; \lambda)$ of \eqref{eq:loss} is
bounded by $B = 1$ and is monotone non-increasing in $\lambda$, and the claims are
exchangeable by assumption. The threshold $\widehat{\lambda}_t$ in
\eqref{eq:sac-threshold} is exactly the conformal-risk-control threshold
\eqref{eq:crc-threshold} computed from the $n_t$ losses of this subsequence.
Applying the conformal risk control theorem \citep{angelopoulos2024crc} to the
subsequence gives $\mathbb{E}[\ell(c; \widehat{\lambda}_t) \mid t(c) = t] \le
\alpha_t$. Since $\mathcal{T}$ partitions the claim space, the tiers are handled by
disjoint calibration subsets and the bounds hold jointly.
\end{proof}

\paragraph{Expected versus per-split risk.}
The guarantee in Proposition~\ref{prop:sac} bounds the \emph{expectation} over the
random draw of calibration and test claims. It does not assert that the realized
tier risk on any single calibration split falls below $\alpha_t$; the empirical
risk fluctuates around its mean from split to split. This distinction drives our
reporting. On held-out claims we measure the realized tier risk
\begin{equation}
\widehat{R}^{\,t}_{\text{test}} \;=\;
\frac{\sum_{c \,:\, t(c)=t} \mathbf{1}[s(c) \ge \widehat{\lambda}_t]\,
\mathbf{1}[y(c)=1]}{\big|\{c : t(c) = t\}\big|},
\label{eq:realized-risk}
\end{equation}
and the violation rate as the fraction of calibration/test splits on which
$\widehat{R}^{\,t}_{\text{test}} > \alpha_t$. An expectation-level guarantee is
consistent with a mean that hugs $\alpha_t$ and a violation rate near one half; we
report both rather than the mean alone, so the reader can see where the guarantee
does and does not bite.

\paragraph{Exchangeability and clustering.}
Proposition~\ref{prop:sac} assumes exchangeability at the claim level. Claims are
in fact nested within questions, and claims from the same answer are correlated, so
claim-level exchangeability is an idealization. Rather than assume it away, we
quantify the resulting finite-sample uncertainty with a question-level cluster
bootstrap (Section~\ref{sec:setup-eval}), resampling whole questions so that the
reported confidence intervals reflect the true unit of correlation.

\subsection{Worked example}
\label{sec:example}
The mechanism is easiest to see on a single answer. Suppose a model is asked how to
take metformin and its answer decomposes into the four scored, tiered claims of
Table~\ref{tab:worked-example}.

\begin{table*}[t]
\centering
\small
\begin{tabular}{llcc}
\toprule
Claim & Tier & $s(c)$ & $y(c)$ \\
\midrule
``Take 5000\,mg per day.'' & \textsc{dangerous} & $0.82$ & $1$ (hallucination) \\
``Avoid metformin in severe kidney disease.'' & \textsc{dangerous} & $0.90$ & $0$ (true) \\
``Metformin derives from French lilac.'' & \textsc{benign} & $0.60$ & $0$ (true) \\
``Metformin was first described in 1922.'' & \textsc{benign} & $0.55$ & $0$ (true) \\
\bottomrule
\end{tabular}
\caption{A single metformin answer decomposed into four scored, tiered claims,
used in the worked example of Section~\ref{sec:sac}.}
\label{tab:worked-example}
\end{table*}

Suppose calibration at the marginal budget $\alpha = 0.10$ returns a single global
threshold $\widehat{\lambda} = 0.70$. Both dangerous claims clear it and are kept,
including the overdose ``5000\,mg'' at $s = 0.82$, while the two true benign claims
are flagged. The harm-blind filter passes the dangerous hallucination because its
score clears a bar set low enough that the abundant, lower-scored benign claims
already hold the pooled rate within budget.

Now apply severity-aware CRC with $\alpha_{\textsc{dangerous}} = 0.05$ and
$\alpha_{\textsc{benign}} = 0.15$. The dangerous tier earns a stricter threshold,
say $\widehat{\lambda}_{\textsc{dangerous}} = 0.88$, and the benign tier a looser
one, $\widehat{\lambda}_{\textsc{benign}} = 0.50$. The overdose claim
($0.82 < 0.88$) is now flagged, the true contraindication ($0.90 \ge 0.88$) is
kept, and both true benign facts ($0.60, 0.55 \ge 0.50$) are surfaced. The scores
and claims are identical; only the allocation of strictness changed. Severity-aware
control catches the dangerous hallucination and retains more true benign
information at the same time, because it stopped spending one budget on two kinds
of error.
